\section{A criterion for the entire relation space}\label{sec:universality}

Fix a prime power $q$. Write
\[
 g_j=\prod_{i=1}^{j}(1-q^{-i}),\qquad
 g_0=1,\qquad g_\infty=\prod_{i=1}^{\infty}(1-q^{-i})>0.
\]
The Gaussian coefficient $\qbinom{d}{s}$ denotes the number of $s$-dimensional
subspaces of $\F_q^d$ and is zero if $s<0$ or $s>d$. In particular,
\begin{equation}\label{eq:gaussian}
 \qbinom{d}{s}=q^{s(d-s)}\frac{g_d}{g_sg_{d-s}}.
\end{equation}
All limits below keep $q$ and every explicitly fixed integer constant.

Let $E_n$ be an $\F_q$-vector space of dimension $D_n\to\infty$, and let
$\nu_n$ be a probability measure on its dual $E_n^*$. For a fixed nontrivial
additive character $\psi$ of $\F_q$, define
\[
 \widehat\nu_n(A)=\mathbb E_{\ell\sim\nu_n}\psi(\ell(A)),\qquad A\in E_n.
\]
For independent $\ell_1,\ldots,\ell_t\sim\nu_n$, the relation space is
$K_n=\bigcap_{i=1}^{t}\ker\ell_i$. Let $K_n^{\rm unif}$ denote the same
intersection when the functionals are independent and uniform on $E_n^*$.
Total variation here is on the finite set of \emph{all subspaces of the same
space $E_n$}; it retains the position of the subspace, not only its dimension.

\begin{theorem}[Fourier criterion for relation-space universality]\label{thm:abstract}
Suppose that there is a constant $\theta<1$ with
\begin{equation}\label{eq:global-bias}
 \sup_{A\ne0}|\widehat\nu_n(A)|\le\theta+o(1).
\end{equation}
Suppose also that for every $\varepsilon>0$ there are sets
$B_n\subset E_n$, containing zero, such that
\begin{equation}\label{eq:exceptional-bias}
 |B_n|\le q^{\varepsilon D_n+o(D_n)},\qquad
 \sup_{A\notin B_n}|\widehat\nu_n(A)|=o(D_n^{-1}).
\end{equation}
For every fixed $u\in\mathbb Z$, at $t=D_n+u$ one has
\[
 \TV\bigl(\mathcal L(K_n),\mathcal L(K_n^{\rm unif})\bigr)\longrightarrow0.
\]
Moreover, for every fixed $s\ge0$,
\begin{equation}\label{eq:moment-limit}
 \mathbb E\qbinom{\dim K_n}{s}
 \longrightarrow\frac{q^{-s(s+u)}}{g_s}.
\end{equation}
\end{theorem}

\begin{proof}
We give the subspace counting and the inversion argument separately, since
convergence of a few dimension moments alone would not prove the stated conclusion.
Write $D=D_n$. For an $s$-dimensional subspace $S\le E_n$, orthogonality gives
\begin{equation}\label{eq:ps}
 p_n(S):=\Pr(\ell|_S=0)
 =q^{-s}\sum_{A\in S}\widehat\nu_n(A),\qquad
 \Pr(S\le K_n)=p_n(S)^t.
\end{equation}
For each fixed $s$ we first prove the stronger summed estimate
\begin{equation}\label{eq:delta-s}
 \Delta_{n,s}:=\sum_{\dim S=s}\left|p_n(S)^t-q^{-st}\right|\longrightarrow0.
\end{equation}
There is nothing to prove for $s=0$. For $s\ge1$, choose $\varepsilon>0$
so small that, for every $1\le r\le s$,
\begin{equation}\label{eq:epsilon-choice}
 r(1-\varepsilon)>\log_q\bigl(1+(q^r-1)\theta\bigr).
\end{equation}
This is possible because $\theta<1$. Choose $B_n$ as in
\eqref{eq:exceptional-bias}. Classify $S$ by
$r=\dim\langle S\cap B_n\rangle$. If $r=0$, then
\[
 p_n(S)=q^{-s}(1+o(D^{-1})),\qquad
 p_n(S)^t=q^{-st}(1+o(1))
\]
uniformly over these subspaces. Their total contribution to
\eqref{eq:delta-s} tends to zero, since
\begin{equation}\label{eq:baseline-moment}
 \qbinom{D}{s}q^{-st}
 =q^{-s(s+u)}\frac{g_D}{g_sg_{D-s}}
 \longrightarrow\frac{q^{-s(s+u)}}{g_s}.
\end{equation}

If $r\ge1$, select a basis of $\langle S\cap B_n\rangle$ from $B_n$ and
extend it to a basis of $S$. Thus the number of such $S$ is at most
\[
 |B_n|^r q^{(s-r)D}\le q^{(s-r+r\varepsilon)D+o(D)}.
\]
Every element of $S$ outside that $r$-space is outside $B_n$. The triangle
inequality in \eqref{eq:ps} therefore gives
\[
 p_n(S)\le q^{-s}\bigl(1+(q^r-1)\theta+o(1)\bigr).
\]
After raising to $t=D+u$ and summing, the exponent of $q$, apart from $o(D)$,
is
\[
 D\bigl[-r(1-\varepsilon)+
       \log_q(1+(q^r-1)\theta)\bigr],
\]
which is strictly negative by \eqref{eq:epsilon-choice}. The sum of the
comparison terms $q^{-st}$ on these subspaces also tends to zero. This
proves \eqref{eq:delta-s}; summing the inclusion probabilities proves
\eqref{eq:moment-limit}.

We now recover exact subspace probabilities. Put
$\mu_j=(-1)^jq^{j(j-1)/2}$. For a fixed $k$-dimensional subspace $T$ and
integer $L\ge0$, let
\[
 F_{n,L}(T)=\sum_{j=0}^{L}\mu_j
       \sum_{\substack{S\supset T\\\dim S=k+j}}\Pr(S\le K_n).
\]
For $m\ge1$, the finite Gaussian identity
\begin{equation}\label{eq:partial-mobius}
 \sum_{j=0}^{L}\mu_j\qbinom mj
 =(-1)^Lq^{L(L+1)/2}\qbinom{m-1}{L}
\end{equation}
follows by induction from the Gaussian recurrence; it holds with the
zero convention also when $L\ge m$. For $m=0$ the sum is one.
Apply this identity with $m=\dim K_n-k$ on the event $T\le K_n$.
The error $F_{n,L}(T)-\Pr(K_n=T)$ has the same sign $(-1)^L$ for every $T$.
Consequently, after summing absolute errors over all $k$-spaces,
\begin{align}
 &\sum_{\dim T=k}|F_{n,L}(T)-\Pr(K_n=T)|\notag\\
 &\quad\le q^{L(L+1)/2}\qbinom{k+L+1}{k}
       \mathbb E\qbinom{\dim K_n}{k+L+1}.
       \label{eq:mobius-tail}
\end{align}
Indeed, when $m\ge L+1$ one has
$\qbinom{m-1}{L}\le\qbinom m{L+1}$, and
\[
 \qbinom dk\qbinom{d-k}{L+1}
 =\qbinom{k+L+1}{k}\qbinom d{k+L+1}.
\]

For fixed $k,L$, \eqref{eq:delta-s} bounds the difference of the two
truncated inversions, summed over $T$, by
\[
 \sum_{j=0}^{L}|\mu_j|\qbinom{k+j}{k}\Delta_{n,k+j}=o(1).
\]
By \eqref{eq:moment-limit}, the limit of the right side of
\eqref{eq:mobius-tail} is at most
$C_q q^{-L^2/2+O_{k,u}(L)}$, and the identical bound holds for the uniform
model. First let $n\to\infty$ with $L$ fixed, then let $L\to\infty$.
The total absolute difference of exact probabilities over $k$-spaces
therefore tends to zero for every fixed $k$.

Finally, the case $s=1$ of \eqref{eq:moment-limit} gives
$\mathbb E q^{\dim K_n}\to1+q^{-u}$. Thus the dimensions are tight,
uniformly for large $n$, and the same is true in the uniform model.
Truncating the dimension and summing the fixed-$k$ conclusions proves
the total variation assertion.
\end{proof}

\begin{corollary}[Universal bounded-offset nullity law]\label{cor:nullity}
Under the hypotheses of \cref{thm:abstract}, $\dim K_n$ converges in total
variation to the distribution
\begin{equation}\label{eq:pi}
 \pi_{q,u}(k)=
 \frac{g_\infty q^{-k(k+u)}}{g_k g_{k+u}},
 \qquad k\ge\max(0,-u).
\end{equation}
The mass is zero at all other integers. In particular, the probability
that the $t$ constraints have the largest possible rank tends to
\begin{equation}\label{eq:full-rank}
 \frac{g_\infty}{g_{|u|}}=\prod_{j>|u|}(1-q^{-j}).
\end{equation}
\end{corollary}

\begin{proof}
For a uniform $(D+u)\times D$ matrix, counting matrices of rank $D-k$
gives
\[
 \Pr(\dim K_n^{\rm unif}=k)
 =q^{-k(k+u)}\frac{g_Dg_{D+u}}{g_kg_{k+u}g_{D-k}}.
\]
Its limit is \eqref{eq:pi}. Tightness, proved above, shows that these
limits sum to one and gives total variation convergence. Set
$k=\max(0,-u)$ for \eqref{eq:full-rank}.
This last matrix-counting law is classical; the new assertion is its
transfer to the structured constraints below.
\end{proof}

\section{Quadratic evaluations}\label{sec:quadratics}

Let $\mathcal Q_n$ denote the space of homogeneous quadratic polynomials
in $n$ variables over $\F_q$, with basis $X_i^2$ and $X_iX_j$, $i<j$.
Its dimension is $n(n+1)/2$, including in characteristic two. Evaluation
on $\F_q^n$ is injective on this space: evaluation at $e_i$ determines
the square coefficients, and evaluation at $e_i+e_j$ determines each
cross coefficient.

\begin{theorem}[Universality for spaces of quadratics]\label{thm:quadratics}
Let $E_n\le\mathcal Q_n$ be arbitrary subspaces with
$D_n:=\dim E_n\to\infty$ and
\begin{equation}\label{eq:quadratic-size}
 \frac{D_n}{n\log D_n}\longrightarrow\infty.
\end{equation}
For independent uniform points
$x_1,\ldots,x_{D_n+u}\in\F_q^n$, define
\[
 K_n=\{f\in E_n:f(x_j)=0\text{ for every }j\}.
\]
For every fixed integer $u$, the entire subspace $K_n$ has total
variation distance tending to zero from the kernel of a uniform
$(D_n+u)\times D_n$ matrix, under any fixed identification of $E_n$
with $\F_q^{D_n}$. Hence \eqref{eq:pi} and \eqref{eq:full-rank} hold.
The estimates are uniform over the choice of $E_n$ at fixed $n,D_n$.
In particular, any bound $D_n\ge\delta n^2$, with fixed $\delta>0$,
suffices.
\end{theorem}

\begin{proof}
For $f\in\mathcal Q_n$, let
$B_f(x,y)=f(x+y)-f(x)-f(y)$ be its polar bilinear form and put
$r(f)=\operatorname{rank}B_f$. With
$b(f)=\mathbb E_x\psi(f(x))$, differencing gives
\begin{equation}\label{eq:quadratic-bias}
 |b(f)|^2
 =\mathbb E_{h}\psi(f(h))\mathbb E_x\psi(B_f(x,h)),
 \qquad |b(f)|\le q^{-r(f)/2}.
\end{equation}
The inner expectation is one on $\operatorname{rad}B_f$ and zero off it;
on that radical, $\psi\circ f$ is an additive character. This proves the
bound (in fact the squared magnitude is either zero or $q^{-r(f)}$).
In odd characteristic, $r(f)=0$ implies $f=0$. In characteristic two,
$r(f)=0$ means $f=\sum_i a_iX_i^2$; if $f\ne0$, its values are uniform
on $\F_q$ by varying a coordinate with $a_i\ne0$, and $b(f)=0$.
Thus for every nonzero $f$,
\[
 |b(f)|\le q^{-1/2}<1.
\]

We need only a coarse count of low-rank forms. The number of bilinear
forms of rank at most $R$ that are symmetric or alternating is at most
\[
 \qbinom n{n-R}q^{R(R+1)/2}
 \le g_\infty^{-1}q^{nR}.
\]
Here a radical subspace of dimension $n-R$ is selected, and a form on
the quotient is selected; this overcounts and is an upper bound.
In odd characteristic the polar map is injective. In characteristic
two its kernel has dimension $n$. Therefore, in either characteristic,
\begin{equation}\label{eq:lowrank-quadratics}
 \#\{f\in\mathcal Q_n:r(f)\le R\}
 \le g_\infty^{-1}q^{n(R+1)}.
\end{equation}
For a given $\varepsilon>0$, which may be reduced to lie in $(0,1)$,
take $R=\lfloor\varepsilon D_n/(2n)\rfloor$ and
$B_n=\{f\in E_n:r(f)\le R\}$. Equations
\eqref{eq:quadratic-bias}--\eqref{eq:lowrank-quadratics} imply
\[
 |B_n|\le q^{\varepsilon D_n/2+o(D_n)},\qquad
 \sup_{f\notin B_n}|b(f)|\le q^{-(R+1)/2}=o(D_n^{-1}).
\]
Here \eqref{eq:quadratic-size} gives $R/\log D_n\to\infty$ and
$n=o(D_n)$. All estimates are uniform in $E_n$.
Apply \cref{thm:abstract}.
\end{proof}

\section{Constraints from random two-planes}\label{sec:planes}

Let $V=\F_q^n$, with $n$ even, and $E_n=\operatorname{Alt}(V)$, the
$D=\binom n2$-dimensional space of alternating bilinear forms. Alternating
means $A(v,v)=0$ for every $v$, including in characteristic two. A uniform
two-plane can be sampled by a uniform ordered independent pair $(x,y)$;
the equation $A(x,y)=0$ depends only on the plane it spans.

\begin{theorem}[The solution space of random plane constraints]\label{thm:planes}
For $t=D+u$ independent uniform two-planes $L_1,\ldots,L_t\le V$, let
\[
 K_n=\{A\in\operatorname{Alt}(V): A|_{L_j\times L_j}=0
          \text{ for every }j\}.
\]
For every fixed $u\in\mathbb Z$,
\[
 \TV\bigl(\mathcal L(K_n),\mathcal L(K_n^{\rm unif})\bigr)\to0,
\]
where $K_n^{\rm unif}$ is the kernel of independent uniform linear
equations on $\operatorname{Alt}(V)$. Consequently its dimension has
limit \eqref{eq:pi}; conditional on any fixed dimension of positive
limiting probability, its position is asymptotically uniform.
\end{theorem}

\begin{proof}
For $A$ of rank $r$, summing over all ordered pairs first gives
$\sum_{x,y}\psi(A(x,y))=q^{2n-r}$. There are
$(q+1)q^n-q$ linearly dependent ordered pairs, and $A(x,y)=0$ on all of
them. The Fourier coefficient for an independent pair is therefore
\begin{equation}\label{eq:plane-fourier}
 \widehat\nu_n(A)=
 \frac{q^{2n-r}-(q+1)q^n+q}{(q^n-1)(q^n-q)}.
\end{equation}
Every nonzero alternating form has rank at least two, so the global
bound is $q^{-2}+o(1)$. The number with rank at most $R$ is at most
\[
 \qbinom n{n-R}q^{\binom R2}
 \le g_\infty^{-1}q^{nR}.
\]
For $R=\lfloor\eta n\rfloor$ and sufficiently small fixed $\eta>0$,
these form an exceptional set as in \eqref{eq:exceptional-bias}; off
that set \eqref{eq:plane-fourier} is
$O_q(q^{-R}+q^{1-n})=o(D^{-1})$. Apply \cref{thm:abstract}.
\end{proof}
